\section{Experimental Protocol}
\label{sec:experiments}

We separate implementation claims from empirical adoption claims.  The model
and losses in \cref{sec:model,sec:consistency} define the method; choosing them
as the production default requires the fixed, multi-seed comparison below.
This separation prevents a low consistency residual from being reported as
evidence that either branch became more accurate.

\subsection{Data partitions and controlled variables}

The primary training source is the target-court singleton synthetic dataset
described in \cref{sec:render_supervision}.  Whole trajectory groups are split
80/10/10, so neighbouring renders never cross partitions.  The validation and
test manifests must report view azimuth, elevation, field of view, coverage
mode, target court, renderer-support flag, and trajectory identifier.  This
metadata enables stratified analysis instead of averaging away difficult
partial views.

A held-out manually annotated real-image set evaluates \KP, line, and
segmentation transfer.  It cannot evaluate translation, rotation, or focal
length unless independent calibration authority is provided; we do not create
pseudo-pose labels from the detector under test.  Synthetic pose metrics and
real dense metrics are therefore reported in separate columns rather than
collapsed into one score.

Across an ablation family we fix the synthetic manifests, \DINO{} checkpoint
and train mode, pose-safe augmentations, optimizer, batch construction, direct
loss weights, image scale (approximately 256 pixels on the short side), and a
15-epoch budget.  Each final comparison uses at least three seeds.  Hyperparameter
search for temperature, Huber threshold, warmup, or auxiliary weight is limited
to a short seed-0 pilot whose selected values are frozen before the multi-seed
run.

\subsection{Ordered architecture and objective ablations}

Architecture is selected before the geometric auxiliary loss is compared.  We
first sweep task-encoder depth with a fixed decoder, then sweep decoder family
and capacity with the chosen encoder.  Only then do we hold the architecture
fixed and compare supervision.  This order avoids attributing a capacity change
to the consistency objective.

\begin{table}[t]
  \centering
  \caption{Ordered ablation plan.  Every final row is repeated with at least
  three seeds after pilot hyperparameters are frozen.}
  \label{tab:ablation_protocol}
  \small
  \begin{tabularx}{\linewidth}{@{}p{0.14\linewidth}p{0.24\linewidth}p{0.25\linewidth}X@{}}
    \toprule
    Stage & Frozen & Varied & Required output \\
    \midrule
    Encoder & Backbone, decoder, direct losses & Task-Transformer depth/width &
      KP/pose accuracy, variance, parameters, time \\
    Decoder & Selected encoder, supervision & linear / progressive / DPT,
      width, taps & Dense accuracy--capacity--speed Pareto \\
    Supervision & Selected architecture, all direct losses & auxiliary off,
      both-grad, two stop-gradient directions & Every GT metric, consistency,
      invalid depth, gradients, memory/time \\
    Robustness & Selected supervision & coverage/pose strata and local-context
      diagnostic & Mean and worst-stratum synthetic/real transfer \\
    \bottomrule
  \end{tabularx}
\end{table}



The core supervision comparison is:
\begin{enumerate}
  \item \textbf{direct-all}: direct \KP, line, segmentation, translation,
    rotation, and focal losses;
  \item \textbf{joint-both}: direct-all plus \cref{eq:cheirality}, with
    gradients to both branches;
  \item \textbf{joint-stopgrad-pose}: the pose projection teaches only the
    dense branch; and
  \item \textbf{joint-stopgrad-dense}: the dense expectation teaches only the
    pose branch.
\end{enumerate}
All four retain identical direct supervision.  The stop-gradient variants
diagnose which branch benefits; they are not alternative test-time graphs.
The auxiliary path is removed at inference, so deployment latency and parameter
count are identical for the four variants.

\subsection{Metrics and non-circular oracles}

Dense keypoints are scored by mean and median Euclidean distance in the model
pixel frame, with per-channel and supervision-eligibility-stratified breakdowns.  Line uses
Dice; segmentation uses mean intersection over union.  Pose uses camera-centre
L2 error in metres, \(\SOthree\) geodesic error in degrees, focal relative
error, and absolute log-focal error.

Two projection metrics must not be confused.  \emph{Pose reprojection error}
projects ground-truth canonical points using only the predicted pose and focal,
then compares with ground-truth image points.  It is an independent pose
metric.  \emph{Branch consistency distance} compares the dense expectation to
the predicted-pose projection.  It is a useful diagnostic but is not a ground-
truth accuracy measure; two jointly wrong branches can make it small.
We additionally record supervision-eligible count, invalid predicted-depth rate, finite
gradient status for both branches, train-step time, and peak accelerator memory.

\subsection{Predeclared adoption rule}

Compared with direct-all, joint-both must improve either GT keypoint mean
distance or pose-only reprojection error by at least 5\% in the three-seed mean.
The other metric and translation/rotation/focal errors may not worsen by more
than 5\% relative.  Absolute degradation must be no larger than 0.01 for line
Dice and segmentation mIoU.  No run may produce a non-finite loss, non-finite
gradient, or all-invalid batch, and step time and peak memory overhead must each
remain within 10\%.  An improvement in consistency distance alone is an
explicit rejection condition.

\subsection{Shortcut-focused diagnostics}

The test manifest is stratified by full, near-full, and partial court coverage;
camera height; azimuth; target-court distance; and number of
supervision-eligible points.  We report the worst stratum alongside the mean.  We also compare full
images with a diagnostic view that retains fixed-size neighbourhoods around
visible junctions while suppressing the remaining context.  This diagnostic is
not used for training or model selection.  If global-pose performance remains
unchanged under the local-only input, the claimed mechanism has not been
empirically demonstrated even if the average dense metric is good.
